%==============================================================================
%  CAPÍTULO — ALEMANIA, FRANCIA E INGLATERRA
%==============================================================================
\chapter{Alemania, Francia e Inglaterra}
\label{cap:comparado-europa-continental}

\fichaCapitulo{Los tres grandes modelos europeos y sus filosofías opuestas.}{La
\emph{Insolvenzordnung} alemana, el \emph{droit des entreprises en difficulté} francés y la
tradición inglesa de la \emph{administration} y los \emph{schemes}: tres respuestas distintas al
mismo problema.}

%==============================================================================
\bloqueTeorico
%==============================================================================

\subsection{Tres filosofías para un mismo problema}

Los tres grandes ordenamientos europeos comparten hoy el marco mínimo de la Directiva
\parencite{ue-directiva-2019-1023}, pero cada uno lo asienta sobre una tradición de fondo distinta
que conviene identificar, porque explica sus soluciones y ofrece a Chile alternativas de diseño.

\begin{description}[leftmargin=0pt,style=nextline,font=\sffamily\bfseries]
  \item[Alemania] Prioriza la \destacar{satisfacción óptima de los acreedores} mediante un
  procedimiento único, técnico y neutral, en que reorganizar o liquidar es una decisión económica y
  no un valor en sí.
  \item[Francia] Prioriza la \destacar{conservación de la empresa y del empleo} como objetivo de
  orden público, con fuerte intervención judicial y una batería de procedimientos preventivos.
  \item[Inglaterra] Prioriza la \destacar{flexibilidad y el control de los acreedores garantizados},
  con instrumentos de mercado, escasa judicialización y gran protagonismo de los profesionales
  privados.
\end{description}

\begin{foco}[El eje que ordena la comparación europea]
Los tres modelos pueden situarse en un eje según a quién sirven primordialmente. Francia, en un
extremo, sirve a la empresa y al empleo; Inglaterra, en el otro, a los acreedores garantizados y a
la eficiencia; Alemania, en el centro, a la satisfacción colectiva de los acreedores. Chile ---con
su procedimiento neutral orientado a la mayor recuperación, pero con superpreferencia laboral y
sustracción de los trabajadores del acuerdo--- se acerca más al modelo alemán con inflexiones
francesas.
\end{foco}

%==============================================================================
\bloqueNormativo
%==============================================================================

\subsection{Alemania: la \emph{Insolvenzordnung}}

La \emph{Insolvenzordnung} (InsO) de 1999 sustituyó el antiguo sistema dual por un
\destacar{procedimiento único de insolvencia} (\emph{Insolvenzverfahren}) dentro del cual se decide,
según convenga a los acreedores, entre la liquidación y la reorganización mediante un plan
(\emph{Insolvenzplan}). La misma puerta de entrada; el destino lo fija el interés de la masa.

Sus rasgos de mayor interés comparado:

\paragraph{Los presupuestos de apertura.} La InsO distingue con precisión técnica tres causales:
la \destacar{insolvencia actual} (\emph{Zahlungsunfähigkeit}, incapacidad de pago), la
\destacar{insolvencia inminente} (\emph{drohende Zahlungsunfähigkeit}, que habilita la apertura
anticipada a instancia del deudor) y el \destacar{sobreendeudamiento} (\emph{Überschuldung}, pasivo
superior al activo). Esta tipología ---más fina que la chilena--- permite graduar el acceso según la
gravedad de la crisis.

\paragraph{El \emph{Schutzschirmverfahren} o «procedimiento del paraguas protector».} Introducido
por la reforma \textsc{esug} de 2012, permite al deudor que aún no es insolvente actual solicitar un
período de hasta tres meses, bajo supervisión de un administrador provisional de su elección, para
elaborar un plan de reestructuración al abrigo de una suspensión de ejecuciones. Es un antecedente
directo del \emph{debtor in possession} atemperado y de la reestructuración preventiva europea.

\paragraph{El \emph{Insolvenzplan}.} El plan agrupa a los acreedores en grupos según su posición
jurídica y económica; cada grupo vota, y opera un mecanismo de \destacar{arrastre entre grupos}
(\emph{Obstruktionsverbot}, prohibición de obstrucción) que permite superar la oposición de un grupo
disidente si no resulta perjudicado respecto de la alternativa liquidatoria y participa
equitativamente del valor. Es el equivalente alemán del \emph{cross-class cram-down}.

\paragraph{La \emph{Restschuldbefreiung}.} La descarga de saldos de la persona natural se obtiene
tras un período de buena conducta que las sucesivas reformas han acortado hasta los
\destacar{tres años} que impone la Directiva. Como en Chile, se condiciona a la probidad del deudor
durante el procedimiento.

\begin{alerta}[La tipología alemana de presupuestos como modelo para Chile]
La \LIR{} chilena no distingue con nitidez entre insolvencia actual, inminente y sobreendeudamiento:
usa causales de hecho tasadas para la liquidación forzosa y no fija umbral para la reorganización. La
tipología alemana ---y en particular la categoría de insolvencia \emph{inminente} como habilitante
de la apertura anticipada--- ofrece un modelo para incentivar la reestructuración temprana, que es
donde se conserva el valor (capítulo \ref{cap:teoria-economica}).
\end{alerta}

\subsection{Francia: el \emph{droit des entreprises en difficulté}}

El derecho francés es el más intervencionista y el que más nítidamente subordina el interés de los
acreedores a la \destacar{conservación de la empresa y del empleo}, jerarquía que el propio Código
de Comercio explicita en el orden de sus objetivos: salvar la empresa, mantener el empleo y, sólo
después, pagar a los acreedores.

Su rasgo distintivo es una \destacar{gradación de procedimientos} según la profundidad de la crisis:

\begin{description}[leftmargin=0pt,style=nextline,font=\sffamily\bfseries]
  \item[Mandat ad hoc y conciliation] Procedimientos \destacar{preventivos y confidenciales}, ante
  un mandatario o conciliador designado por el presidente del tribunal, para negociar con los
  principales acreedores antes de toda cesación de pagos. La confidencialidad ---ausente en el
  concurso chileno, esencialmente público--- busca evitar el estigma que precipita la crisis.
  \item[Sauvegarde] Procedimiento de \destacar{salvaguarda} abierto a instancia del deudor que aún
  no está en cesación de pagos, con suspensión de ejecuciones y elaboración de un plan. Es el
  equivalente funcional de la reorganización preventiva.
  \item[Redressement judiciaire] Reorganización judicial del deudor ya en cesación de pagos, con
  posible desplazamiento de la administración.
  \item[Liquidation judiciaire] Liquidación cuando la reorganización es manifiestamente imposible.
\end{description}

\paragraph{Las clases de partes afectadas.} La transposición de la Directiva introdujo en 2021 las
\emph{classes de parties affectées} y el arrastre entre clases, morigerando el tradicional poder del
juez francés de imponer el plan.

\paragraph{El superprivilegio de los salarios.} Francia lleva la protección del crédito laboral más
lejos que casi cualquier ordenamiento, mediante el \emph{superprivilège des salaires} y un régimen
de garantía (\textsc{ags}) que \destacar{adelanta a los trabajadores el pago} de sus créditos con
cargo a un fondo, subrogándose éste luego en el concurso. La comparación es aleccionadora para
Chile, cuya superpreferencia laboral no cuenta con un fondo de garantía que anticipe el pago.

\begin{foco}[La confidencialidad preventiva: lo que Chile no tiene]
El aporte francés más difícil de replicar y quizá el más valioso es la fase
\destacar{preventiva y confidencial}. El concurso chileno es público desde su inicio ---publicación
en el \boletin{}---, lo que dispara el estigma, la fuga de proveedores y el retiro de líneas de
crédito justo cuando la empresa aún podría salvarse. Un instrumento de negociación confidencial
previo, al modo del \emph{mandat ad hoc}, permitiría intentar la reestructuración sin desencadenar
la crisis que se quiere evitar.
\end{foco}

\subsection{Inglaterra: \emph{administration}, CVA y \emph{schemes}}

El derecho inglés ---regido por la \emph{Insolvency Act} 1986 y la \emph{Companies Act} 2006, con la
reforma de la \emph{Corporate Insolvency and Governance Act} 2020--- ofrece el modelo más flexible,
menos judicializado y más orientado al mercado. Sus instrumentos centrales:

\paragraph{La \emph{administration}.} Un \emph{insolvency practitioner} ---profesional privado
licenciado--- asume el control de la compañía con el objetivo, jerarquizado por la ley, de rescatarla
como empresa en marcha; si no es posible, obtener para los acreedores un resultado mejor que la
liquidación; y, en último término, realizar los activos. Notablemente, la \emph{administration} puede
iniciarse \destacar{sin intervención judicial}, por nombramiento del propio deudor o del acreedor
titular de una \emph{floating charge}.

\paragraph{El \emph{Company Voluntary Arrangement} (CVA).} Acuerdo entre la compañía y sus
acreedores, aprobado por mayoría del 75\,\% en valor, que vincula a los disidentes. Es el equivalente
funcional del acuerdo de reorganización, con una tramitación notablemente más ligera.

\paragraph{El \emph{scheme of arrangement} y el \emph{restructuring plan}.} El \emph{scheme}, figura
de la \emph{Companies Act}, permite reestructurar el pasivo mediante un acuerdo sancionado por el
tribunal y aprobado por clases. La reforma de 2020 añadió el \emph{restructuring plan}, que incorpora
un \destacar{\emph{cross-class cram-down}} de estilo estadounidense: el tribunal puede imponer el
plan a clases disidentes si ninguna queda peor que en la alternativa relevante y una clase «in the
money» lo aprueba. El \emph{scheme} inglés ha atraído reestructuraciones de compañías de todo el
mundo ---el llamado \emph{forum shopping} hacia Londres---.

\paragraph{La \emph{floating charge} y el \emph{administrator}.} La garantía flotante inglesa
---gravamen sobre un conjunto de activos cambiantes que se «cristaliza» al abrirse la insolvencia---
otorga a su titular un poder de control del procedimiento sin equivalente en el derecho chileno de
garantías, más rígido y basado en la prenda e hipoteca sobre bienes determinados.

\begin{alerta}[La desjudicialización inglesa y sus límites para Chile]
El modelo inglés confía la insolvencia a profesionales privados con mínima intervención judicial, lo
que gana en celeridad y flexibilidad. Pero descansa sobre una infraestructura ausente en Chile: un
cuerpo consolidado de \emph{insolvency practitioners}, una cultura de garantías flotantes y una
judicatura mercantil especializada de respaldo. Importar los instrumentos sin ese sustrato
---como advierte la experiencia de otros trasplantes--- produce forma sin función.
\end{alerta}

%==============================================================================
\bloquePractico
%==============================================================================

\subsection{Cuadro comparativo de los modelos europeos}

\begin{table}[htbp]
\centering
\footnotesize
\caption{Alemania, Francia e Inglaterra frente a la Ley \Nro 20.720}
\label{tab:cap28-europa}
\begin{tabularx}{\textwidth}{@{}>{\raggedright\arraybackslash}p{0.20\textwidth}
                                 >{\raggedright\arraybackslash}p{0.19\textwidth}
                                 >{\raggedright\arraybackslash}p{0.19\textwidth} X@{}}
\toprule
\textbf{Rasgo} & \textbf{Alemania} & \textbf{Francia} & \textbf{Inglaterra} \\
\midrule
Prioridad & Satisfacción de acreedores & Empresa y empleo & Flexibilidad y garantizados \\
\addlinespace
Entrada & Procedimiento único & Gradación de procedimientos & Instrumentos múltiples \\
\addlinespace
Fase confidencial & No & Sí (mandat ad hoc) & Parcial (schemes) \\
\addlinespace
Arrastre de clases & Sí (Obstruktionsverbot) & Sí (desde 2021) & Sí (plan de 2020) \\
\addlinespace
Judicialización & Media & Alta & Baja \\
\addlinespace
Crédito laboral & Preferente & Superprivilegio + fondo AGS & Preferente limitado \\
\addlinespace
Descarga persona natural & 3 años & Sí & Sí \\
\bottomrule
\end{tabularx}
\end{table}

\begin{foco}[Qué puede tomar Chile de cada modelo europeo]
De \destacar{Alemania}, la tipología fina de presupuestos ---en especial la insolvencia inminente
como llave de la reestructuración temprana--- y el arrastre entre grupos del \emph{Insolvenzplan}.
De \destacar{Francia}, la fase preventiva confidencial que evita el estigma, y el fondo de garantía
que anticipa el pago de los salarios. De \destacar{Inglaterra}, la agilidad del acuerdo vinculante
por mayoría y el \emph{cross-class cram-down} del \emph{restructuring plan}. Tres reformas de
distinta ambición, pero todas convergentes en el mismo diagnóstico chileno: reestructurar antes,
con menos estigma y con capacidad de superar minorías obstruccionistas.
\end{foco}
